\BourbakiAppendixExercises{Exercises}

\begin{enumerate}
\def\labelenumi{\arabic{enumi})}
\tightlist
\item
  Let K be a commutative field and E a vector space over K. We denote by \(\operatorname{End}_{\mathrm{K}}^{\mathrm{pf}}(\mathrm{E})\) the subset of \(\operatorname{End}_{\mathrm{K}}(\mathrm{E})\) consisting of the endomorphisms of which a power belongs to \(\operatorname{End}_{\mathrm{K}}^{\mathrm{f}}(\mathrm{E})\) . If \(u \in \operatorname{End}_{\mathrm{K}}^{\mathrm{pf}}(\mathrm{E})\) , then we set \(I_{u} = \bigcap_{n \geqslant 0} Im u^{n}\) and \(\operatorname{Tr} u = \operatorname{Tr}(u | I_{u})\) ; this notion coincides with that defined in No.~2 for an endomorphism of finite rank. a) Let \(u \in \operatorname{End}_{\mathrm{K}}^{\mathrm{pf}}(\mathrm{E})\) ; let V be a subspace of E stable under u, and let \(u_{V}\) and \(u_{E/V}\) be the endomorphisms of V and E/V induced by u. Prove the equality \(\operatorname{Tr} u = \operatorname{Tr} u_{V} + \operatorname{Tr} u_{E/V}\) .
\end{enumerate}

\begin{enumerate}
\def\labelenumi{\alph{enumi})}
\setcounter{enumi}{1}
\item
  Let U be a subspace of \(\operatorname{End}_{\mathbf{K}}(\operatorname{E})\) . Suppose that there exists an integer n such that every product of n elements of U has finite rank. Prove that the mapping \(Tr: U \to K\) is linear.
\item
  Let F be a vector space over K, and let \(u: E \to F\) and \(v: F \to E\) be homomorphisms such that \(u \circ v \in \operatorname{End}_{\mathrm{K}}^{\mathrm{pf}}(\mathrm{F})\) . We have \(v \circ u \in \operatorname{End}_{\mathrm{K}}^{\mathrm{pf}}(\mathrm{E})\) and \(\operatorname{Tr} v \circ u = \operatorname{Tr} u \circ v\) .
\end{enumerate}

¶ 2) Keep the notation of the previous exercise. If V and W are subspaces of E, then we denote by V \(\prec\) W the relation ``W has finite codimension in V + W.''

\begin{enumerate}
\def\labelenumi{\alph{enumi})}
\item
  Let V be a subspace of E. We denote by \(\mathfrak{A}\) (resp. \(\mathfrak{a},\mathfrak{b}\) ) the set of endomorphisms \(u\) of V such that \(u(\mathrm{V})\prec \mathrm{V}\) (resp. \(u(\mathrm{E})\prec \mathrm{V}\) , \(u(\mathrm{V})\prec 0\) ). Prove that \(\mathfrak{A}\) is a K-subalgebra of \(\operatorname{End}_{\mathrm{K}}(\mathrm{E})\) , that \(\mathfrak{a}\) and \(\mathfrak{b}\) are two-sided ideals of \(\mathfrak{A}\) , with sum \(\mathfrak{A}\) , and that we have \(\mathfrak{a} \cap \mathfrak{b} \subset \operatorname{End}_{\mathrm{K}}^{\mathrm{pf}}(\mathrm{E})\) .
\item
  Let \(u, u'\) be permutable elements of \(\mathfrak{A}\) ; we can write \(u = a + b\) and \(u' = a' + b'\) with \(a, a'\) in \(\mathfrak{a}\) and \(b, b'\) in \(\mathfrak{b}\) . Prove that the element \([a, a'] = aa' - a'a\) of \(\mathfrak{A}\) belongs to \(\mathfrak{a} \cap \mathfrak{b}\) and that its trace does not depend on the chosen decompositions of \(u\) and \(u'\) ; we denote it by \(\langle u, u' \rangle\) .
\item
  Let A be a commutative subalgebra of \(\mathfrak{A}\) . Prove that there exists a unique K-linear mapping \(\operatorname{Res}:\Omega_{\mathrm{K}}(\mathrm{A})\to \mathrm{K}\) such that \(\operatorname{Res}(udv) = \langle u,v\rangle\) for all \(u,v\) in A (use the identity \([u,vw] + [v,wu] + [w,uv] = 0\) ).
\item
  Let \(u, v \in A\) ; suppose \(v(V) \subset V\) . Let \(\pi\) be a projector of E with image V. Prove the equality \(\operatorname{Res}(u dv) = \operatorname{Tr}[\pi u, v]\) . Deduce that if \(u(V) \subset V\) , then we have \(\operatorname{Res} u dv = 0\) (observe that the endomorphism \([\pi u, v]\) has square zero).
\item
  Let \(u \in \mathrm{A}\) , \(n \in \mathbf{N}\) . Prove that we have \(\operatorname{Res} u^n du = 0\) . If \(u\) is invertible in \(\mathrm{A}\) , then we have \(\operatorname{Res} u^{-n} du = 0\) for \(n \geqslant 2\) ; if moreover \(u(\mathrm{V}) \subset \mathrm{V}\) , then we have \(\operatorname{Res} u^{-1} du = \dim_{\mathrm{K}} \mathrm{V} / u(\mathrm{V})\) .
\end{enumerate}

\(f)\) Take \(\mathrm{E} = \mathrm{K}((\mathrm{X}))\) (IV, §4, No.~9, p.~38), \(\mathrm{V} = \mathrm{K}[[\mathrm{X}]]\) , and let \(\mathrm{A} = \mathrm{E}\) act on itself by homotheties. Prove that the assumption of \(c)\) is satisfied. Let \(f(\mathrm{X}) = \sum_{n} a_{n} \mathrm{X}^{n} d\mathrm{X}\) be an element of \(\mathrm{K}((\mathrm{X}))\) . Prove that \(\operatorname{Res} f(\mathrm{X}) d\mathrm{X}\) is equal to \(a_{-1}\) . Deduce that if \(g\) is an element of \(\mathrm{K}[[\mathrm{X}]]\) such that \(g(0) = 0\) and \(g'(0) \neq 0\) , then the coefficients of \(\mathrm{X}^{-1}\) in the formal series expansions of \(f(\mathrm{X})\) and \(f(g(\mathrm{X})) g'( \mathrm{X})\) are the same.

\backmatter
